% GENERATED FILE. Edit canonical lessons, then run npm run generate:books.
% canonical-source-hash: fnv1a64:76397ab865a85223
% canonical-lessons: UR-C154-tika, UR-C154-dharkan, UR-C154-pasina, UR-C154-ulti, UR-C154-chakkar

\chapter{Vaccination, Heartbeat, Sweat, Vomiting, Dizziness}
\label{ch:ur-vaccination-heartbeat-sweat-vomiting-dizziness}
\hlchaptermodality{154}

\begin{chapteropening}
\textbf{By the end of this chapter:} \emph{I can say a vaccination, a heartbeat, sweat, vomiting and dizziness.}

Complete the last lesson of chapter 154: I can say a vaccination, a heartbeat, sweat, vomiting and dizziness.
\end{chapteropening}

\section[\texorpdfstring{\ur{ٹیکا} (ṭīkā)}{ٹیکا (ṭīkā)}]{\ur{ٹیکا} (ṭīkā) — a vaccination}

\label{lesson:UR-C154-tika}



\begin{quote}
Before the new one: say the Urdu for an eyelid, then the Urdu for an eyebrow.
\end{quote}

\subsection*{You'll want to know: \ur{ٹیکا}}
\textbf{\ur{ٹیکا}} — \emph{ṭīkā} — \textquotedblleft{}a vaccination\textquotedblright{}.

\begin{cousinweb}
\textbf{\ur{ٹیکا}} (\emph{ṭīkā}) is a tadbhava, a word inherited and reshaped by sound change. It comes from \emph{*ṭikka}, \textquotedblleft{}mark, spot\textquotedblright{}, a form CDIAL reconstructs from its descendants (CDIAL 5458). CDIAL glosses the Hindi word \textquotedblleft{}caste-mark\textquotedblright{}; Urdu also uses it for a vaccination.
\end{cousinweb}

\begin{grammarlens}[title={What you've built}]
One more word for this chapter.
\end{grammarlens}

\subsection*{Your turn}
\begin{itemize}
\raggedright
  \item \emph{Say it:} \emph{ṭīkā}
  \item \emph{Say it:} \emph{ṭīkā}, once more
  \item \emph{Say it:} say \emph{ṭīkā}
  \item \emph{Recall:} say the Urdu for an eyelid, then the Urdu for an eyebrow, then say \emph{ṭīkā} again
\end{itemize}

\begin{tcolorbox}[breakable,colback=teal!4,colframe=teal!35!black,title={Before you move on}]
What does \ur{ٹیکا} mean? (\textquotedblleft{}A vaccination\textquotedblright{}.) Say it once more.
\end{tcolorbox}
\section[\texorpdfstring{\ur{دھڑکن} (dhaṛkan)}{دھڑکن (dhaṛkan)}]{\ur{دھڑکن} (dhaṛkan) — a heartbeat}

\label{lesson:UR-C154-dharkan}



\begin{quote}
Before the new one: say the Urdu for an eyebrow, then the Urdu for a vaccination.
\end{quote}

\subsection*{You'll want to know: \ur{دھڑکن}}
\textbf{\ur{دھڑکن}} — \emph{dhaṛkan} — \textquotedblleft{}a heartbeat\textquotedblright{}.

\begin{cousinweb}
\textbf{\ur{دھڑکن}} (\emph{dhaṛkan}) is the noun of \emph{dhaṛaknā}, \textquotedblleft{}to throb\textquotedblright{}, which CDIAL lists under a reconstructed \emph{*dhaḍ}, \textquotedblleft{}throb, beat\textquotedblright{} (CDIAL 6711).
\end{cousinweb}

\begin{grammarlens}[title={What you've built}]
One more word for this chapter.
\end{grammarlens}

\subsection*{Your turn}
\begin{itemize}
\raggedright
  \item \emph{Say it:} \emph{dhaṛkan}
  \item \emph{Say it:} \emph{dhaṛkan}, once more
  \item \emph{Say it:} say \emph{dhaṛkan}
  \item \emph{Recall:} say the Urdu for an eyebrow, then the Urdu for a vaccination, then say \emph{dhaṛkan} again
\end{itemize}

\begin{tcolorbox}[breakable,colback=teal!4,colframe=teal!35!black,title={Before you move on}]
What does \ur{دھڑکن} mean? (\textquotedblleft{}A heartbeat\textquotedblright{}.) Say it once more.
\end{tcolorbox}
\section[\texorpdfstring{\ur{پسینہ} (pasīnā)}{پسینہ (pasīnā)}]{\ur{پسینہ} (pasīnā) — sweat}

\label{lesson:UR-C154-pasina}



\begin{quote}
Before the new one: say the Urdu for a vaccination, then the Urdu for a heartbeat.
\end{quote}

\subsection*{You'll want to know: \ur{پسینہ}}
\textbf{\ur{پسینہ}} — \emph{pasīnā} — \textquotedblleft{}sweat\textquotedblright{}.

\begin{cousinweb}
\textbf{\ur{پسینہ}} (\emph{pasīnā}) is a tadbhava, a word inherited and reshaped by sound change. It comes from \emph{*prasvidana}, \textquotedblleft{}sweating\textquotedblright{}, a form CDIAL reconstructs from its descendants (CDIAL 8895). CDIAL also allows \emph{prasvinna}, \textquotedblleft{}sweated\textquotedblright{}.
\end{cousinweb}

\begin{grammarlens}[title={What you've built}]
One more word for this chapter.
\end{grammarlens}

\subsection*{Your turn}
\begin{itemize}
\raggedright
  \item \emph{Say it:} \emph{pasīnā}
  \item \emph{Say it:} \emph{pasīnā}, once more
  \item \emph{Say it:} say \emph{pasīnā}
  \item \emph{Recall:} say the Urdu for a vaccination, then the Urdu for a heartbeat, then say \emph{pasīnā} again
\end{itemize}

\begin{tcolorbox}[breakable,colback=teal!4,colframe=teal!35!black,title={Before you move on}]
What does \ur{پسینہ} mean? (\textquotedblleft{}Sweat\textquotedblright{}.) Say it once more.
\end{tcolorbox}
\section[\texorpdfstring{\ur{الٹی} (ulṭī)}{الٹی (ulṭī)}]{\ur{الٹی} (ulṭī) — vomiting}

\label{lesson:UR-C154-ulti}



\begin{quote}
Before the new one: say the Urdu for a heartbeat, then the Urdu for sweat.
\end{quote}

\subsection*{You'll want to know: \ur{الٹی}}
\textbf{\ur{الٹی}} — \emph{ulṭī} — \textquotedblleft{}vomiting\textquotedblright{}.

\begin{cousinweb}
\textbf{\ur{الٹی}} (\emph{ulṭī}) is \textquotedblleft{}an overturning\textquotedblright{}, from \emph{ulaṭnā}, \textquotedblleft{}to be overturned\textquotedblright{}, which CDIAL lists under a reconstructed \emph{*ullaṭati}, \textquotedblleft{}turns over\textquotedblright{}, in its form \emph{*ullaṭyate}, \textquotedblleft{}is upset\textquotedblright{} (CDIAL 2368).
\end{cousinweb}

\begin{grammarlens}[title={What you've built}]
One more word for this chapter.
\end{grammarlens}

\subsection*{Your turn}
\begin{itemize}
\raggedright
  \item \emph{Say it:} \emph{ulṭī}
  \item \emph{Say it:} \emph{ulṭī}, once more
  \item \emph{Say it:} say \emph{ulṭī}
  \item \emph{Recall:} say the Urdu for a heartbeat, then the Urdu for sweat, then say \emph{ulṭī} again
\end{itemize}

\begin{tcolorbox}[breakable,colback=teal!4,colframe=teal!35!black,title={Before you move on}]
What does \ur{الٹی} mean? (\textquotedblleft{}Vomiting\textquotedblright{}.) Say it once more.
\end{tcolorbox}
\section[\texorpdfstring{\ur{چکر} (chakkar)}{چکر (chakkar)}]{\ur{چکر} (chakkar) — dizziness; a round}

\label{lesson:UR-C154-chakkar}



\begin{quote}
Before the new one: say the Urdu for sweat, then the Urdu for vomiting.
\end{quote}

\subsection*{You'll want to know: \ur{چکر}}
\textbf{\ur{چکر}} — \emph{chakkar} — \textquotedblleft{}dizziness; a round\textquotedblright{}.

\begin{cousinweb}
\textbf{\ur{چکر}} (\emph{chakkar}) goes back to Sanskrit \emph{cakra}, \textquotedblleft{}wheel, circle\textquotedblright{}. CDIAL lists the inherited Hindi \emph{cāk}, \textquotedblleft{}wheel\textquotedblright{}, under it (CDIAL 4538), but not \emph{chakkar} itself, which keeps more of the Sanskrit shape.
\end{cousinweb}

\begin{grammarlens}[title={What you've built}]
One more word for this chapter.
\end{grammarlens}

\subsection*{Your turn}
\begin{itemize}
\raggedright
  \item \emph{Say it:} \emph{chakkar}
  \item \emph{Say it:} \emph{chakkar}, once more
  \item \emph{Say it:} say \emph{chakkar}
  \item \emph{Recall:} say the Urdu for sweat, then the Urdu for vomiting, then say \emph{chakkar} again
\end{itemize}

\begin{tcolorbox}[breakable,colback=teal!4,colframe=teal!35!black,title={Before you move on}]
What does \ur{چکر} mean? (\textquotedblleft{}Dizziness; a round\textquotedblright{}.) Say it once more.
\end{tcolorbox}
